\section*{الفصل \ref{ch:b3:heterocycles} --- كيمياء الحلقات غير المتجانسة}

\begin{solution}{exo:b3:heterocycles:1}
الأوكسازول: يعطي O اثنين (زوج حر في النظام $\pi$، وآخر في المستوي)، وN واحدًا، 
وذرات الكربون الثلاث ثلاثة: 6. الثيازول: الشيء نفسه مع S: 6. البيريميدين: واحد من كل
ذرة من الذرات الست، وزوجا النيتروجين الحران كلاهما في المستوي: 6. البيورين: تسع ذرات،
يعطي نيتروجين N--H اثنين والأخرى واحدًا لكل منها: 10، وهو عدد هوكل
($4n + 2$، $n = 2$).
\end{solution}

\begin{solution}{exo:b3:heterocycles:2}
البيبيريدين (11.1) $>$ DMAP (9.6) $>$ الإيميدازول (7.0) $>$ البيريدين (5.2) $\gg$ البيرول (نحو $-4$،
مبرتنًا على الكربون).
\end{solution}

\begin{solution}{exo:b3:heterocycles:3}
2-برومو بيرول (والبيرول شديد التفاعلية إلى حد أنه يتبرمن تبرمنًا متعددًا بسهولة)؛
و2-برومو فيوران؛ و3-برومو إندول.
\end{solution}

\begin{solution}{exo:b3:heterocycles:4}
2،5-ثنائي ميثيل فيوران؛ 1،2،5-ثلاثي ميثيل بيرول؛ 2،5-ثنائي ميثيل ثيوفين.
\end{solution}

\begin{solution}{exo:b3:heterocycles:5}
2-ميثوكسي بيريدين. فإضافة الميثوكسيد عند C2 تضع الشحنة السالبة \omterm{def:b3:heterocycles:snar}{لمعقد
مايزنهايمر} على النيتروجين؛ أما عند C3 فلا يمكن نشرها إلا على ذرات الكربون، فيكون
الوسيط، والحالة الانتقالية التي تسبقه، أعلى طاقة بكثير.
\end{solution}

\begin{solution}{exo:b3:heterocycles:6}
2-أمينو بيريدين (في صورة ملحه الصوديومي، ثم مبرتنًا عند المعالجة). يُضاف أيون الأميد
عند C2، فيعطي ناتج إضافة $\sigma$ أنيونيًا فيه H و$\ce{NH2}$ على C2 والشحنة
السالبة على نيتروجين الحلقة؛ ثم يفقد هيدريدًا، يأخذ بروتونًا من
مجموعة الأمينو ويغادر في صورة $\ce{H2}$.
\end{solution}

\begin{solution}{exo:b3:heterocycles:7}
يعيد أكسجين N-أكسيد زوجًا حرًا إلى الحلقة: فتحمل البنى الرنينية
شحنة سالبة عند C2 وC4 وC6. وتهاجم المحبات للإلكترونات C4 (لأن C2
قريب من النيتروجين المشحون إيجابًا). ثم يُنزع الأكسجين بثلاثي
كلوريد الفوسفور، فيعطي 4-نتروبيريدين.
\end{solution}

\begin{solution}{exo:b3:heterocycles:8}
2-هيدروكسي بيريدين (مجموعة OH على C2) و2-بيريدون (N--H وC=O). يحتفظ البيريدون بستة
إلكترونات $\pi$ في الحلقة: فبنيته الرنينية الثنائية الأيون، ذات $\mathrm{O^-}$ 
و$\mathrm{N^+}$ الذي يتقاسم زوجه الحر، أولات بيريدينيوم؛ وهو أميد يكمل
الزوج الحر لنيتروجينه السداسية، والروابط الهيدروجينية القوية للمجموعتين C=O وN--H
تفضّله في الماء.
\end{solution}

\begin{solution}{exo:b3:heterocycles:9}
2-نتروبنزالدهيد، ومكافئان من أسيتو أسيتات الميثيل، والأمونيا: وهو
حاصر قنوات الكالسيوم النيفيديبين.
\end{solution}

\begin{solution}{exo:b3:heterocycles:10}
الإين-هيدرازين المتجه نحو مجموعة $\ce{CH2}$ (الداخلي، الأكثر استبدالًا) يعطي
2،3-ثنائي ميثيل إندول؛ والمتجه نحو مجموعة الميثيل (الطرفي) يعطي 2-إيثيل إندول.
والأحماض الضعيفة تعطي أساسًا الأول، عبر الإين-هيدرازين الأكثر استقرارًا؛ 
والأحماض الأقوى ترفع نسبة 2-إيثيل إندول.
\end{solution}

\begin{solution}{exo:b3:heterocycles:11}
في $\mathrm S_\mathrm NAr$ تكون الخطوة المحددة للسرعة هي الإضافة، التي لا تكسر الرابطة C--X؛
والفلور الأكثر كهرسلبية يثبّت الحالة الانتقالية الأنيونية
\omterm{def:b3:heterocycles:snar}{ومعقد مايزنهايمر} أكثر. أما في $\mathrm S_\mathrm N2$ فتنكسر الرابطة C--X في
الحالة الانتقالية، وتغادر الرابطة C--I الضعيفة أسرع.
\end{solution}

\begin{solution}{exo:b3:heterocycles:12}
البيريدين: ثلاث إشارات بالنسبة 2\,:\,1\,:\,2، أولاها قرب \qty{8.6}{ppm}. البيرول: إشارتان
متساويتان عند 6.7 و\qty{6.2}{ppm} وإشارة N--H عريضة قرب \qty{8.1}{ppm}. الفيوران: إشارتان متساويتان عند
7.4 و\qty{6.4}{ppm}، دون N--H. الثيوفين: إشارتان متساويتان عند 7.33 و\qty{7.12}{ppm}، متقاربتان،
بين إشارتي الفيوران.
\end{solution}

\begin{solution}{pb:b3:heterocycles:1}
\textbf{1.} $\ce{C6H5NHNH2 + (CH3)2CO -> C6H5NHN=C(CH3)2 + H2O}$.

\textbf{2.} يُضاف نيتروجين الأمين إلى كربون الكربونيل، ثم يُفقد الماء، كما في
تكوّن الإيمين؛ والحمض ينشّط مجموعة الكربونيل ويحوّل مجموعة OH في ناتج الإضافة
إلى مجموعة مغادرة.

\textbf{3.} مجموعة $\ce{NH2}$ الطرفية: فالزوج الحر للنيتروجين الآخر مترافق مع
حلقة الفينيل وأكثر إعاقة.

\textbf{4.} $\qty{108.0}{g/mol}$; $\qty{10.0}{g}/\qty{108.0}{g/mol} = \qty{0.0926}{mol}$.

\textbf{5.} الهيدرازون $\ce{C9H12N2}$، \qty{148.0}{g/mol}: $\qty{0.0926}{mol} \times \qty{148.0}{g/mol} = \qty{13.7}{g}$.

\textbf{6.} $\ce{C6H5-NH-NH-C(CH3)=CH2}$.

\textbf{7.} كربون الحلقة في الموضع \textit{ortho}، وكربون الحلقة المرتبط بالنيتروجين، وذرتا النيتروجين، وكربون الهيدرازون
ومجموعة $\ce{CH2}$ الطرفية.

\textbf{8.} تنكسر الرابطة N--N؛ وتتكوّن رابطة C--C بين كربون الموضع \textit{ortho} ومجموعة $\ce{CH2}$.

\textbf{9.} ستة إلكترونات، وكل المكوّنات أحادية الوجه: مسموحة حراريًا.

\textbf{10.} بوصفه حمض لويس يحفز التصاوغ، ويرتبط بذرة نيتروجين 
ويُضعف الرابطة N--N، ويساعد الأمونيا على المغادرة.

\textbf{11.} صار كربون الموضع \textit{ortho} من نوع $sp^3$، حاملًا هيدروجينًا: وفقد ذلك البروتون
يستعيد حلقة البنزين، وهذه قوة دافعة كبيرة.

\textbf{12.} تهاجم مجموعة $\ce{NH2}$ الأنيلينية المتكوّنة كربون الإيمين، فتغلق
الحلقة الخماسية (2-أمينو إندولين).

\textbf{13.} الأمونيا؛ وتكوّن حلقة البيرول العطرية في الإندول هو الذي يدفعها.

\textbf{14.} $\ce{C9H12N2 -> C9H9N + NH3}$.

\textbf{15.} \qty{131.0}{g/mol}.

\textbf{16.} $\qty{0.0926}{mol} \times \qty{131.0}{g/mol} = \qty{12.1}{g}$.

\textbf{17.} عند C3، الموضع الأكثر محبة للنواة في الإندول، والحر هنا؛ والهجوم هناك
يحتفظ بحلقة البنزين سليمة في الوسيط.

\textbf{18.} $\ce{C6H5NHNH-C(=CH2)-CH2CH3}$ (الطرفي) و$\ce{C6H5NHNH-C(CH3)=CH-CH3}$ (الداخلي).

\textbf{19.} 2-إيثيل إندول و2،3-ثنائي ميثيل إندول.

\textbf{20.} الداخلي أكثر استبدالًا؛ والأحماض القوية ترفع حصة
الطرفي.

\textbf{21.} 2،3-ثنائي ميثيل إندول، عبر الإين-هيدرازين الأكثر استقرارًا والأكثر
استبدالًا.

\textbf{22.} يتكاثف الأسيتالدهيد مع نفسه ويعطي طريق الإين-هيدرازين مردودات
ضعيفة؛ فيُصنع الإندول بدلًا من ذلك من هيدرازون حمض البيروفيك، فيعطي
حمض إندول-2-كربوكسيليك، الذي يُنزع منه الكربوكسيل.

\textbf{23.} يجب فصل مزيج من المتماكبات الموضعية، وهذا يكلّف من المردود؛ فيُفضَّل
كيتون متناظر أو طريق يثبّت الانتقائية الموضعية.

\textbf{24.} يمكن أن يعطي 10.0 g من الفينيل هيدرازين \qty{12.1}{g} من 2-ميثيل إندول على الأكثر.
\end{solution}
